\section{Directional source stability forces the same canonical point orders}

Fix an actual formal source ideal $K\subset\mathbb C[[x_1,\ldots,x_n]]$,
a point chart $(e,b)$, its genuine linear source shear $S$, and the literal
direction $v_e=1$, $v_i=b_i$ for $i\ne e$. Write $J$ for the independently
defined whole formal strict ideal. The natural source qualification below is
that $\delta_vP\in\operatorname{in}K$ for every polynomial $P$ whose whole
formal image belongs to $\operatorname{in}K$. It is a derivation-stability
condition, not an order, diagram-comparison, shadow or HS certificate.

\begin{theorem}[Directional sparsity of the same canonical source forms]
\label{hir:actual-canonical-directional-form-sparsity}
\uses{hir:actual-directional-normal-polynomial-dehom-euler,hir:actual-reduced-formal-vertex-members,hir:actual-homogeneous-initial-ideal-diagram-equality,hir:actual-formal-normal-representative-unique}
\lean{Hironaka.HironakaActualDirectionalPivotNormality.canonical_source_form_directional_zero,Hironaka.HironakaActualDirectionalPivotNormality.canonical_source_dehomogeneous_homogeneous}
\leanok
Assume the stated source directional stability. For every SAME canonical
reduced source series $G_a$ of vertex degree $w=|a|$, its literal degree-$w$
source polynomial $P_a$ satisfies $\delta_vP_a=0$; consequently
$\operatorname{dehom}_{e,b}P_a$ is genuinely homogeneous of degree $w$.
\end{theorem}
\begin{proof}
The source form $P_a$ belongs to the whole source initial ideal, so its
directional derivative does also. At any exponent $\beta\in\Gamma(K)$,
upward closure puts $\beta+\mathbf1_i$ in $\Gamma(K)$. Its coefficient
in $G_a$ is zero unless this exponent is $a$. That exception is impossible:
minimality of the vertex $a$ and $\beta\le a$ would imply $a\le\beta$,
contradicting the positive $i$-increment. Every partial-derivative summand
therefore has zero coefficient on the ENTIRE source diagram. The initial
ideal has exactly this diagram, so genuine normal-representative uniqueness
makes the whole directional derivative zero. The previously proved actual
Euler coefficient formula now reads
$(w-|\beta|)[X^\beta]\operatorname{dehom}P_a=0$. A nonzero coefficient
forces $|\beta|=w$, since a nonzero integer is nonzero in $\mathbb C$.
This is precisely whole homogeneous degree $w$, including the zero form.
\end{proof}

\begin{theorem}[The entire point shadow is the sheared source initial ideal]
\label{hir:actual-directional-whole-shadow-source-initial}
\uses{hir:actual-canonical-directional-form-sparsity,hir:actual-reduced-canonical-homogeneous-initial-span,hir:actual-point-homogeneous-dehom-shear-identity,hir:actual-point-exceptional-restriction-dehomogenized-polynomial}
\lean{Hironaka.HironakaActualDirectionalPivotNormality.pivot_projection_strict_eq_sheared_source_initial,Hironaka.HironakaActualDirectionalPivotNormality.sheared_source_initial_component_mem}
\leanok
Under the same natural qualification,
$\rho_e(J)=S(\operatorname{in}K)$ as ENTIRE ideals. Every homogeneous
component of a member of this whole sheared ideal remains in it.
\end{theorem}
\begin{proof}
The genuine strict-family generation theorem uses exactly the original
canonical $G_a$ with their actual source point powers. For each such
normalized member, the genuine exceptional congruence identifies its entire
pivot restriction with $\operatorname{dehom}P_a$. The latter is pivot free
by its literal polynomial substitution. The preceding theorem makes it
homogeneous of degree $w$, and the genuine shear/dehom identity gives
$\operatorname{dehom}P_a=S(P_a)$. Mapping the entire original canonical
initial-form span and the entire strict-family span proves both ideal
inclusions, hence the asserted equality. For component closure, express an
arbitrary whole ideal member as an actual finite sum $\sum c_a S(P_a)$.
Its degree-$q$ component is the sum of $[c_a]_{q-w}S(P_a)$ for $w\le q$;
terms with $w>q$ are zero. Every term lies in the same whole ideal.
No homogeneous closure or projection equality is assumed.
\end{proof}

\begin{theorem}[All current exceptional colons force every same canonical order]
\label{hir:actual-directional-canonical-order-of-current-colons}
\uses{hir:actual-directional-whole-shadow-source-initial,hir:actual-current-colons-homogeneous-initial-pivot-cancellation,hir:actual-canonical-higher-band-dehom-initial-independence,hir:actual-point-normalized-all-pivot-coefficient-bands,hir:actual-current-point-factor-source-band-dictionary}
\lean{Hironaka.HironakaActualDirectionalPivotNormality.canonical_directional_quotient_power_of_colons}
\leanok
Assume source directional stability and the actual whole colons
$(J+\mathfrak m^{q+1}):e=J+\mathfrak m^q$ for every $q$.
For every SAME canonical $G_a$ and every genuine quotient
$e^{|a|}H=\pi_{e,b}^*G_a$, one has $H\in\mathfrak m^{|a|}$.
\end{theorem}
\begin{proof}
Otherwise take the genuine first nonzero current degree $q<|a|$ of $H$.
Its degree-$q$ form $F$ lies in the whole current initial ideal. We prove
that any factorization $F=e^rL$, with $L$ homogeneous of degree $q-r$ and
in that same initial ideal, has $L=0$, by induction on $q-r$.
The actual degree-lift theorem and the preceding component closure put
$\rho_e(L)$ in $S(\operatorname{in}K)$. The literal coefficient-band
formula identifies it with degree $q-r$ of
$\operatorname{dehom}([G_a]_{|a|+r})$. The actual point power
$H\in\mathfrak m^q$ supplies all lower bands zero. When $r=0$, the first
source-form theorem makes the degree-$q$ band zero because $q<|a|$.
When $r>0$, the source band has degree strictly greater than $|a|$;
the SAME canonical reduced family makes it normal on the ENTIRE original
diagram, and genuine directional first-band independence makes its band
zero. Thus $\rho_e(L)=0$. If its degree is zero, this forces $L=0$.
Otherwise the actual pivot kernel gives $L=eL'$; the genuine initial
pivot cancellation derived from the given colons retains $L'$ inside the
same current initial ideal. Its degree is one less and the factorization
becomes $F=e^{r+1}L'$, so induction applies. At $r=0$ this contradicts the
nonzero genuine initial form of $H$. Every source band and quotient in
this argument belongs to the original prescribed family.
\end{proof}

\begin{theorem}[Qualified original arbitrary-direction persistence implies every canonical order]
\label{hir:actual-original-directional-persistent-canonical-orders-qualified}
\uses{hir:actual-directional-canonical-order-of-current-colons,hir:actual-original-point-persistent-hs-full-initial-colons,hir:actual-point-centre-whole-analytic-strict-generators}
\lean{Hironaka.HironakaActualDirectionalPivotNormality.original_directional_persistent_hs_orders_of_source_stable}
\leanok
For an original analytic point ideal $I$, assume the stated natural
source directional stability for its ACTUAL Taylor ideal and actual entire
source/current HS equality in chart $(e,b)$. Then every genuine normalized
quotient of the SAME original canonical formal source series retains its
original vertex order.
\end{theorem}
\begin{proof}
The actual analytic-to-formal strict-ideal dictionary identifies the entire
current Taylor ideal with the independently defined formal strict ideal.
Actual original HS persistence, the real Bennett comparison and the finite
jet rigidity theorem supply every required current exceptional colon.
Apply the preceding theorem to the literal source Taylor ideal, chart and
prescribed quotient. This invokes no new source family and no additional
current-order or standardness hypothesis.
\end{proof}

The directional source qualification is explicit in this chapter; its
derivation from arbitrary-direction original persistence is separate geometry.
This chapter does not assert the original-only arbitrary-direction criterion,
positive-dimensional-centre maximal standardness, or any history termination.
